\par\Needspace{6\baselineskip}\medskip
\hypertarget{native-151-article109}{}\pdfbookmark[2]{109.}{article-109}
\begin{center}109.\end{center}
Si in hyperbola signa $\alpha$, $C$, $c$ in eadem significatione accipimus, ut in art. 99, habemus ex aequationibus VIII, IX art. 21
\[
\begin{aligned}
r'\cos v'-r\cos v&=-\tfrac12\left(c-\frac1c\right)\left(C-\frac1C\right)\alpha\\[.7em]
r'\sin v'-r\sin v&=\tfrac12\left(c-\frac1c\right)\left(C+\frac1C\right)\alpha\BareRoot{(ee-1)},
\end{aligned}
\]
\noindent adeoque
\[
\rho=\tfrac12\alpha\left(c-\frac1c\right)\BareRoot{\left(ee\left(C+\frac1C\right)^2-4\right)}.
\]
\noindent\SourceLeaf{152}Supponamus $\gamma$ esse quantitatem per aequationem $\gamma+\dfrac1\gamma=e\left(C+\dfrac1C\right)$ determinatam: cui quum manifesto \emph{duo} valores sibi invicem reciproci satisfaciant, adoptamus eum qui est maior quam 1. Ita fit
\[
\rho=\tfrac12\alpha\left(c-\frac1c\right)\left(\gamma-\frac1\gamma\right)\mathpunct{\cdot}
\]
\noindent Porro fit
\[
\begin{aligned}
r+r'&=\tfrac12\alpha\left(e\left(c+\frac1c\right)\left(C+\frac1C\right)-4\right)\\[.7em]
&=\tfrac12\alpha\left(\left(c+\frac1c\right)\left(\gamma+\frac1\gamma\right)-4\right),
\end{aligned}
\]
\noindent adeoque
\[
\begin{aligned}
r+r'+\rho&=\alpha\left(\BareRoot{c\gamma}-\BareRoot{\frac1{c\gamma}}\right)^2\\[.7em]
r+r'-\rho&=\alpha\left(\BareRoot{\frac\gamma c}-\BareRoot{\frac c\gamma}\right)^2\mathpunct{\cdot}
\end{aligned}
\]
\noindent Statuendo itaque
\[
\BareRoot{\frac{r+r'+\rho}{4\alpha}}=m,\qquad\BareRoot{\frac{r+r'-\rho}{4\alpha}}=n,
\]
\noindent erit necessario $\BareRoot{c\gamma}-\BareRoot{\dfrac1{c\gamma}}=2m$; ad decidendam vero quaestionem, utrum $\BareRoot{\dfrac\gamma c}-\BareRoot{\dfrac c\gamma}$ fiat $=+2n$ an $=-2n$, inquirere oportet, utrum $\gamma$ maior an minor sit quam $c$: sed ex aequatione 8 art. 99 facile sequitur, casum priorem locum habere, quoties $2f$ sit infra $180^0$, posteriorem, quoties $2f$ sit supra $180^0$. Denique ex eodem art. habemus
\[
\begin{aligned}
\frac{kt}{\alpha^{\frac32}}&=\tfrac12\left(\gamma+\frac1\gamma\right)\left(c-\frac1c\right)-2\log c\\[.8em]
&=\tfrac12\left(c\gamma-\frac1{c\gamma}\right)-\tfrac12\left(\frac\gamma c-\frac c\gamma\right)-\log c\gamma+\log\frac\gamma c\\[.8em]
&=2m\BareRoot{(1+mm)}\mp2n\BareRoot{(1+nn)}\\[.5em]
&\qquad-2\log\bigl(\BareRoot{(1+mm)}+m\bigr)\pm2\log\bigl(\BareRoot{(1+nn)}+n\bigr)
\end{aligned}
\]
\noindent signis inferioribus semper ad casum $2f>180^0$ spectantibus. Iam
\[
\log\bigl(\BareRoot{(1+mm)}+m\bigr)
\]
\noindent facile evolvitur in seriem sequentem
\[
m-\tfrac13\cdot\tfrac12m^3+\tfrac15\cdot\frac{1\mathpunct{.}3}{2\mathpunct{.}4}m^5-\tfrac17\cdot\frac{1\mathpunct{.}3\mathpunct{.}5}{2\mathpunct{.}4\mathpunct{.}6}m^7+\text{etc.}
\]
\noindent Hoc sponte colligitur ex $\mathrm d\log\bigl(\BareRoot{(1+mm)}+m\bigr)=\dfrac{\mathrm dm}{\BareRoot{(1+mm)}}\mathpunct{\cdot}$ Prodit itaque
\[
\begin{aligned}
&2m\BareRoot{(1+mm)}-2\log\bigl(\BareRoot{(1+mm)}+m\bigr)\\[.6em]
&\qquad=4\left(\tfrac13m^3-\tfrac15\cdot\tfrac12m^5+\tfrac17\cdot\frac{1\mathpunct{.}3}{2\mathpunct{.}4}m^7-\text{etc.}\right)
\end{aligned}
\]
\noindent\SourceLeaf{153}et perinde formula alia prorsus similis, si $m$ cum $n$ permutatur. Hinc tandem colligitur, si statuatur
\[
\begin{aligned}
T={}&\tfrac16(r+r'-\rho)^{\frac32}-\tfrac1{80}\cdot\frac1\alpha(r+r'-\rho)^{\frac52}\\[.6em]
&+\tfrac3{1792}\cdot\frac1{\alpha\alpha}(r+r'-\rho)^{\frac72}-\tfrac5{18432}\cdot\frac1{\alpha^3}(r+r'-\rho)^{\frac92}+\text{etc.}
\end{aligned}
\]
\[
\begin{aligned}
U={}&\tfrac16(r+r'+\rho)^{\frac32}-\tfrac1{80}\cdot\frac1\alpha(r+r'+\rho)^{\frac52}\\[.6em]
&+\tfrac3{1792}\cdot\frac1{\alpha\alpha}(r+r'+\rho)^{\frac72}-\tfrac5{18432}\cdot\frac1{\alpha^3}(r+r'+\rho)^{\frac92}+\text{etc.,}
\end{aligned}
\]
\noindent fieri $kt=U\mp T$, quae expressiones cum iis, quae in art. 107 traditae sunt, omnino coincidunt, si illic $a$ in $-\alpha$ mutetur.
\par
Ceterum hae series tum pro ellipsi tum pro hyperbola ad usum practicum tunc imprimis sunt commodae, ubi $a$ vel $\alpha$ valorem permagnum obtinet, i. e. ubi sectio conica magnopere ad parabolae similitudinem vergit. In tali casu etiam ad solutionem problematis supra tractati (artt. 85---105) adhiberi possent: sed quoniam, nostro iudicio, ne tunc quidem brevitatem solutionis supra traditae praebent, huic methodo fusius exponendae non immoramur.
\par\bigskip\begin{center}\rule{.23\linewidth}{.35pt}\end{center}
